% Properly normalized independent bivariate Gaussian; sigma_x=1.2, sigma_y=0.55.
\clearpage
\section*{Visual laboratory: a probability density is a geometric field}
\noindent For independent normal measurements with
\(\sigma_x=1.2,\;\sigma_y=0.55\), the joint density is
\[
p(x,y)=\frac{1}{2\pi(1.2)(0.55)}
\exp\!\left[-\frac12\!\left(\frac{x^2}{1.2^2}+
\frac{y^2}{0.55^2}\right)\right].
\]
Height and color encode the same density. A taller peak does not imply
greater total probability: the integral over the full plane is one.
\begin{figure}[!htbp]
\centering
\begin{tikzpicture}
\begin{groupplot}[
 group style={group size=2 by 1,horizontal sep=10mm},
 booktwoaxis,width=0.44\textwidth,height=63mm,
 xmin=-3,xmax=3,ymin=-2,ymax=2,
 xlabel={$x$},ylabel={$y$},
 point meta min=0,point meta max=0.245,
 colormap name=booktwothermal
]
\nextgroupplot[title={3D probability density},view={40}{28},zmin=0,zmax=0.26]
\addplot3[surf,shader=interp,draw=none,domain=-3:3,y domain=-2:2,samples=29,samples y=23]
{exp(-0.5*((x/1.2)^2+(y/0.55)^2))/(2*pi*1.2*0.55)};
\nextgroupplot[title={Top-down density map},view={0}{90},axis equal image]
\addplot3[surf,shader=interp,draw=none,domain=-3:3,y domain=-2:2,samples=29,samples y=23]
{exp(-0.5*((x/1.2)^2+(y/0.55)^2))/(2*pi*1.2*0.55)};
\end{groupplot}
\end{tikzpicture}
\caption{Analytical bivariate normal density with covariance
\(\Sigma=\operatorname{diag}(1.44,0.3025)\).
The figures depict the formula above, not observed sample frequencies or a confidence region.
Equal-density contours are ellipses; their probability content must be calculated, not inferred from color alone.}
\label{fig:book02-gaussian-uncertainty}
\end{figure}
\noindent\textbf{Reconstruct it:} draw seeded samples, estimate their
covariance, and compare empirical contour coverage with the intended
\(\chi^2_2\) probability levels across many independent replications.
